% frontmatter/preface.tex

\addchap{Предисловие}

Подключить платёжный шлюз можно за день: SDK, сумма, номер карты,
ответ \texttt{00}~\enquote{Approved} -- дело сделано. Это и есть
ловушка.

Ощущение лёгкости рассеивается, когда приходит первый
\textit{chargeback}. Эмитент начинает отклонять каждую пятую
транзакцию с~кодом \texttt{05}~\enquote{Do Not Honor} без объяснения
причин. Процессинг формирует клиринговый файл, в~котором суммы
не сходятся с~авторизациями, и~непонятно, ошибка это или нет.
Аудитор спрашивает, почему PAN лежит в~логах, и~разговор уходит
в~нарушения PCI DSS. К~этому моменту выясняется, что \texttt{00}
ничего не гарантировал.

\medskip

Платёжная индустрия живёт парадоксом: подключить шлюз -- один день;
понять, что произошло после \texttt{Approved}, -- годы. Знания
в~ней передаются устно и~оттого ведут себя предсказуемо плохо.

Они \textit{фрагментарны}: специалист по эмиссии не объяснит
клиринг, специалист по PCI DSS не открывал EMV Book~2, архитектор
не читал положение Банка России №~762-П. Каждый защищает свой
участок и~называет соседний \enquote{не моей зоной ответственности}.

Они \textit{искажаются при передаче}: \enquote{Visa штрафует
за~магнитную полосу} -- верно не всегда и~не в~каждом регионе;
у~правила \enquote{3-D~Secure переносит ответственность} -- десяток
исключений. Без первоисточника отличить правило от его упрощённой
версии невозможно, а~цена ошибки оплачивается чужими деньгами.

Их \textit{не хватает на растущую команду}: одного носителя
знаний на двадцать инженеров недостаточно, и~новый человек
приходит ровно к~тому же фольклору, через который прошёл его
предшественник.

\medskip

Эта книга выросла из заметок, которые я делал для себя несколько
лет, пытаясь свести воедино три источника, говорящих о~платежах
на разных языках: технические спецификации (прежде всего EMVCo
и~PCI~SSC), правила платёжных систем и~российскую нормативную базу.
Каждый написан в~своей традиции, использует свою терминологию
и~предполагает своего читателя; согласовать их оказалось задачей
более трудоёмкой, чем я предполагал вначале.

Я держался одного правила: ни одно утверждение в~книге не должно
опираться на цеховое знание без указания первоисточника. Там, где
источники расходятся -- а~они расходятся чаще, чем хотелось бы, --
расхождение отмечено явно. Там, где я не уверен, я говорю об этом
прямо: \enquote{не знаю}, а~не \enquote{кажется как будто бы}.
Эта калька с~английского \textit{seems like} проросла в~российский
финтех как мягкий способ пофантазировать; в~книге её нет.

Предварительных знаний о~платежах не требуется. Требуется
готовность открыть EMV Book~2 на странице~67, чтобы убедиться,
что я ничего не перепутал; без этого большая часть материала
останется абстракцией.

\hfill\textit{О.Д.}

\hfill\textit{\the\year}
